\section{Resiliência nos adaptadores de saída}
\label{sec:resiliencia-quarkus}

A aula~2 e o capítulo~\ref{cap:ponte-quarkus} apresentaram retry,
timeout, fallback, circuit breaker e bulkhead, e o SmallRye Fault
Tolerance que o Quarkus usa. A pergunta arquitetural é \emph{onde}
colocar as anotações \texttt{@Retry}, \texttt{@Timeout},
\texttt{@Fallback}, \texttt{@CircuitBreaker} e \texttt{@Bulkhead} que o
guia \citetitle{doc:quarkus-fault-tolerance} mostra.

\paragraph{Na borda, perto da chamada que falha.}
A falha de rede é um detalhe do adaptador: o domínio não sabe que a
cobrança é HTTP. As anotações vão no adaptador de saída que implementa a
porta:

\begin{lstlisting}[language=Java]
package br.com.turma.pedidos.adaptador.saida.cobranca;
@ApplicationScoped
class CobrancaHttp implements CobrancaGateway {
    @Inject @RestClient CobrancaApi api;   // package-private, sem reflexão

    @Timeout(800) @Retry(maxRetries = 2) @CircuitBreaker
    public ResultadoCobranca cobrar(Pedido p) { /* converte e chama api */ }
}
\end{lstlisting}

\paragraph{O que a porta devolve quando tudo falha.}
O \texttt{@Fallback} é uma decisão de negócio disfarçada de técnica:
\enquote{cobrança indisponível} vira pedido pendente, erro para o usuário
ou fila? Exponha esse resultado no tipo de retorno da porta
(\texttt{ResultadoCobranca} com um caso \texttt{Indisponivel}, por
exemplo, com um \texttt{sealed interface} da
seção~\ref{sec:sealed-pattern}), para que o caso de uso decida e o teste
unitário cubra a decisão.

\paragraph{Bulkhead e virtual threads.}
O bulkhead limita execuções concorrentes. Com virtual threads, a regra
da seção~\ref{sec:virtual-threads} é a mesma: não limite com pool de
threads, limite o recurso escasso. O bulkhead no adaptador é o lugar
natural para isso.

\paragraph{Testando.}
Os padrões da aula~2 se testam no adaptador: um \texttt{@QuarkusTest}
com o serviço externo simulado (por exemplo, com o Toxiproxy da aula~2)
mostra se timeout e circuit breaker se comportam como configurado. O
caso de uso, por sua vez, testa só a reação a
\texttt{Indisponivel}, com um fake da porta.
